% !PRINTED 56
\weberpage{56}

Ist nun $\varkappa = 1$, so sind nach der Definition von $g$ die Zahlen

\[
1,\ g,\ g^2 \ldots g^{q-2}
\]

nach dem Modul $q$ alle von einander verschieden, während $g^{q-1}$ wieder congruent mit 1 ist.

Ist aber $\varkappa > 1$, so ist:

\weq{4}{g^{q-1} = 1 + hq,}

und nach unserer Voraussetzung über $g$ ist $h$ nicht durch $q$ theilbar. Wenn wir die Gleichung (4) in die Potenz $q$ erheben, und rechts den binomischen Lehrsatz anwenden, so ergiebt sich

\[
g^{q(q-1)} = 1 + hq^2 + h^2\, \frac{q^3(q-1)}{2} + \cdots,
\]

also

\weq{5}{g^{q(q-1)} = 1 + h_1 q^2,}

worin

\[
h_1 = h + h^2\, \frac{q(q-1)}{2} + \cdots
\]

nicht durch $q$ theilbar ist. (Das wäre für $q = 2$ nicht mehr richtig und darum verlangt die Primzahl 2 eine andere Behandlung.)

Erheben wir (5) nochmals in die $q^{\mathrm{te}}$ Potenz, so ergiebt sich

\[
g^{q^2(q-1)} = 1 + h_2 q^3,
\]

und so können wir fortfahren und erhalten für jeden beliebigen positiven Exponenten $\lambda$

\weq{6}{g^{q^{\lambda-1}(q-1)} = 1 + hq^\lambda,}

worin $h$ eine durch $q$ nicht theilbare ganze Zahl ist. Setzen wir zur Abkürzung:

\weq{7}{q^{\varkappa-1}(q-1) = \varphi(q^\varkappa) = c,}

so ist also

\weq{8}{g^c \equiv 1\ (\mathrm{mod}\ q^\varkappa),}

und es ist noch nachzuweisen, dass $c$ der kleinste positive Exponent ist, für den die Congruenz (8) erfüllt ist. Nehmen wir an, es sei $e$ dieser kleinste Exponent, also

\weq{9}{g^e \equiv 1\ (\mathrm{mod}\ q^\varkappa),}

so muss $e$ ein Theiler von $c$ sein, weil sonst die Congruenz (8) auch erfüllt wäre, wenn $c$ durch den Rest der Division von $c$ durch $e$, der kleiner als $e$ ist, ersetzt wird.

Andererseits muss $e$ durch $q-1$ theilbar sein, weil die in (9) enthaltene Congruenz $g^e \equiv 1\ (\mathrm{mod}\ q)$ nur für solche Exponenten, die durch $q-1$ theilbar sind, besteht.
